\chapter{\label{cha:conclusion}Conclusion and future work}

\section{Conclusion}
\label{sec:conclusion}

This thesis studied spectral bias in a controlled Neural Tangent Kernel (NTK)
setting: a one-hidden-layer ReLU network trained on functions on the unit circle
\(S^1\). The continuum infinite-width model with the uniform input measure gives
the reference case. In this limit, the NTK is rotation-invariant, so the
corresponding operator is diagonal in the Fourier basis. Training then has an
explicit frequency-wise form: each Fourier component of the residual decays at
a rate determined by the corresponding kernel eigenvalue.

The first part of the thesis studied what changes after replacing the continuum
measure by finitely many sampled points. Exact Fourier diagonalization is lost,
but the Fourier description remains accurate on fixed low-frequency blocks. The
sampled Fourier Gram matrix \(G\), the matrix \(H\) describing the sampled
operator on the retained Fourier coordinates, and the difference
\(A\Phi-\Phi\Lambda^{(n)}\) all become smaller as the sample size increases.
This explains why the finite-sample corrected eigenvalues predict the dominant
early-time residual dynamics. The preconditioning experiment supports the same
interpretation: using only the diagonal eigenvalue prediction removes most of
the observed rate separation on the retained block.

The second part of the thesis studied finite width in the frozen-kernel setting.
For each fixed Fourier block, the frozen finite-width tangent operator at
initialization has expectation equal to the continuum Fourier block and
concentrates around it as the width increases. The numerical experiments show
the same trend both at the matrix level and at the eigenspace level: the matched
empirical eigenspaces become more aligned with the corresponding Fourier
frequency subspaces as width increases.

The evolving finite-width experiments add a further qualification. The frozen
finite-width theory shows that the tangent kernel at initialization approaches
the continuum Fourier operator as width increases. During training, however, the
kernel can move away from this frozen baseline. Empirically, this movement has
two visible effects. It can reduce alignment with some Fourier subspaces,
especially at higher frequencies, while increasing the average kernel strength
in the lowest Fourier blocks. At small width, this low-frequency strength
increase is associated with faster decay of the dominant target modes and lower
final loss. At larger width, the frozen kernel is already closer to the
continuum reference, and the evolving-kernel advantage decreases.

The main conclusion is that the Fourier basis gives the right reference point,
but its role changes across regimes. It is exact for the continuum
infinite-width operator. It remains approximately valid after finite sampling
on fixed low-frequency blocks. It also describes the frozen finite-width
tangent kernel well when the width is large. For the evolving finite-width
kernel, the Fourier subspaces are best viewed as diagnostics: they reveal how
training changes alignment, residual decay, and the distribution of kernel
strength across frequencies. A rigorous theory matching these evolving-kernel
observations remains open.

\section{Future work}
\label{sec:future-work}

The most direct next step is to develop theory for the evolving finite-width
tangent kernel. The experiments in this thesis compare frozen and evolving
dynamics, but the analysis remains empirical. A natural goal is to control the
time-dependent block \(C_t^{(K)}\), its drift from initialization, and its effect
on the residual dynamics.

A second direction is to prove a joint result that includes finite sampling,
finite width, and training time. The present thesis studies finite sampling and
frozen finite width separately, and studies the evolving finite-width kernel
empirically. In actual training runs, all three effects occur at once.

A third direction is to sharpen the concentration bounds. The current proofs use
entrywise concentration, coarse tail bounds, and union bounds. These arguments
are sufficient to prove the correct decreasing trends, but they are conservative
as numerical estimates. Direct operator-norm bounds, variance-sensitive
estimates, or concentration inequalities for U- and V-statistics may give bounds
closer to the observed errors.

A fourth direction is to study the alignment diagnostics more directly. The
projector-error and principal-angle plots suggest that different Fourier
subspaces can have different finite-width stability. A more detailed analysis
could measure which Fourier blocks are most strongly mixed by the finite-width
operator and how this depends on the eigenvalues and gaps of the continuum
operator.

Finally, the analysis should be extended beyond the uniform circle setting.
Non-uniform input densities, higher-dimensional spheres, other activations, and
different parameterizations all change the relevant kernel operator and its
spectrum. Studying these cases would clarify which conclusions are special to
the explicit Fourier model on \(S^1\), and which reflect a broader mechanism
behind spectral bias.
